\chapter*{Abstract}

La presente relazione descrive il project work sviluppato in collaborazione con TP, avente come oggetto la progettazione e la realizzazione di T.A.L.O.S. \textit{(Taranto 2026 AI Live Operator Support)}, una soluzione basata su intelligenza artificiale multimodale per il supporto al customer care dei Giochi del Mediterraneo 2026 a Taranto.

Il progetto prende avvio da una precisa esigenza operativa: superare la frammentazione delle fonti informative che caratterizza gli eventi internazionali, garantendo così un flusso di dati e di informazioni coerente e accessibile. In un contesto caratterizzato da cittadini, turisti e stakeholder provenienti da Paesi differenti, la disponibilità di un canale informativo unico rappresenta un elemento decisivo per migliorare l’esperienza dell’utente e ridurre il carico sui canali tradizionali di assistenza.

T.A.L.O.S. è stato quindi progettato come un sistema completo, in grado di supportare l’utente secondo le sue esigenze al fine di rendere la sua esperienza quanto più soddisfacente possibile. La soluzione integra un’interfaccia responsive, consultabile da qualsiasi dispositivo, attraverso cui l’utente può accedere a un chatbot basato su una pipeline di \textit{Retrieval-Augmented Generation} (RAG) applicata a una base di conoscenza specifica per il contesto dei Giochi del Mediterraneo 2026. Il sistema supporta inoltre la multimodalità, consentendo l’invio di richieste in formato testuale, vocale e visuale, che vengono elaborate e ricondotte a una rappresentazione comune per generare risposte testuali multilingua in linguaggio naturale.

Quando una richiesta non può essere gestita in modo affidabile in autonomia, oppure quando l’utente esprime insoddisfazione rispetto alla risposta ricevuta, il sistema consente l’apertura di un ticket verso un operatore di customer care. L’operatore dispone di una dashboard riservata per consultare i ticket, visualizzare il contesto della conversazione e gestire la richiesta. È inoltre presente una pagina dedicata all'organizzazione della knowledge base, attraverso cui è possibile inserire nuovi contenuti informativi tramite interfaccia grafica alla portata di tutti, indipendentemente dal livello di competenza tecnica.

Infine, il progetto include un flusso di testing automatizzato basato su un dataset di casi simulati, costruiti per rappresentare possibili interazioni utente. La valutazione misura diverse metriche quantitative relative a diversi aspetti del sistema, riconducibili a specifici \textit{Key Performance Indicator} (KPI).